\documentclass[12pt]{article}

\usepackage{amsmath}
%\usepackage{english}
\usepackage{url}
% \usepackage[dvips]{epsfig}
\usepackage{graphicx}

%\graphicspath{{/home/ritz/Teaching/fcn/fig/}}

%----------------- layout ----------------------------------------

\oddsidemargin=0in
\textwidth=6in
\topmargin=-0.7in
\textheight=9.5in
\baselineskip=16pt
\font\mittel=cmr17
\font\roman=cmr12
\font\klein=cmr9
\font\mgross=cmr10
\newfont{\CAP}{cmcsc10 scaled 900}
\pagestyle{empty}
%parindent=0mm

%----------------- macros ----------------------------------------
\newcommand{\conc}[1]{[{\rm C}]_{\text{#1}}}
\newcommand{\units}[1]{\text{#1}}
\newcommand{\kalium}{$\text{K}^+$}
\newcommand{\natrium}{$\text{Na}^+$}
\newcommand{\chlor}{$\text{Cl}^-$}
\newcommand{\wasser}{$\text{H}_2\text{O}$}
%-----------------------------------------------------------------


\begin{document}

\parbox{2cm}{
\includegraphics[width=1.8cm]{../bccnlogo.pdf}
}
\parbox{11cm}{
\begin{center}
\large HUMBOLDT-UNIVERSIT"AT \hskip 0.1 cm ZU \hskip 0.1 cm BERLIN
\vskip 0.1 true cm
\mgross BERNSTEIN CENTER FOR\\ COMPUTATIONAL NEUROSCIENCE
\end{center}
}
\parbox{2cm}
{
\hfill
\includegraphics[width=1.8cm]{../hublogo.pdf}
}

\vskip 0.6 true cm

\leftline{\CAP Humboldt-Universit"at zu Berlin 
\hfill Phone: 030/2093-9110}
\leftline{\CAP Philippstr. 13 House 6
             \hfill Fax: 030/2093-6771}
\leftline{\CAP \hfill webpage: http://www.bccn-berlin.de/}


\vskip 0.8 true cm
\centerline{\bf Models of Neural Systems, WS 2010/11}
\centerline{\bf Project 4: Action potential propagation}
\centerline{Project presentation and report submission: February,
15th, 2011}

\vskip 0.8 true cm

\newcommand{\Na}{\mathrm{Na}}
\newcommand{\K}{\mathrm{K}}
\newcommand{\Leak}{\mathrm{L}}

\vskip 0.8 true cm

\noindent The aim of this project is to model the propagation of an action
potential along an axon. The relationship between the membrane current
$i_m$ and the voltage along an axon is given by the equation:

\begin{equation}
    c_m\frac{\partial V}{\partial t} =
    \frac{1}{2ar_L}\frac{\partial}{\partial
    x}\left(a^2\frac{\partial V}{\partial x}\right)-i_m+i_e,
\end{equation}
where $a$ is the radius of the axon, $r_L$ is intracellular
resistivity.

The ionic current flowing through a patch of axonal membrane $i_m$ is
well-described by Hodgkin-Huxley model:

\begin{equation}
    i_m=g_\Na m^3h(V-E_\Na)+g_\K n^4(V-E_\K)+g_\Leak(V-E_\Leak),
\end{equation}
where $m$, $n$, $h$ are Hodgkin-Huxley-type activation variables.

Combining these two equations leads to a partial differential equation
which can be computed numerically by a multi-compartmental approximation.
In a nonbranching cable, each compartment is coupled to two neighbours
and the equations for the membrane potentials of the compartments are:

\begin{equation}
   c_m\frac{dV_\mu}{dt}=-i^\mu_m +
    \frac{I^\mu_e}{A_\mu}+g_{\mu,\mu+1}(V_{\mu+1} -V_\mu)+g_{\mu,\mu-1}(V_{
    \mu-1}-V_{\mu}),
\end{equation}
where $\mu$ labels the compartments, $I^\mu_e$ is the total electrode
current flowing into the compartment $\mu$, and $A_\mu$ is its surface
area. The constant $g_{\mu,\mu-1}$ determines the resistive coupling
of the compartments and for nonbranching cables can be shown to be
equal to $g_{\mu,\mu-1}=a/(2r_LL^2)$. This defines a system of
ordinary differential equations which can be solved by generalized Euler
methods.

\vskip 0.8 true cm

\noindent {\bf Problems}

\vskip 0.8 true cm

\begin{enumerate}
    \item Numerically solve the cable equation for passive membrane
        (i.e. $i_m=(V-V_{rest})/r_m$). Compare the solution to the analytical
        solution. Take $r_m=1\mathrm{M\Omega mm^2}$,
        $r_L=1\mathrm{k\Omega mm^2}$.
    \item Implement the Hodgkin-Huxley model of action potential
        propagation. Solve the partial differential equation using the
        Crank-Nicolson method (see e.g. Chapter 6.6B of \cite{dayan2001}).
        Take $r=0.238$ mm and $r_L=35.4$ $\mathrm{\Omega cm}$.
    \item Initiate an action potential on one end of the axon by
        inserting a current in the terminal compartment. 
    \item Determine action potential propagation velocity as a
        function of the axon radius.
    \item Generate action potentials from both ends of the axon. Show
        that they annihilate when they collide.
    \item Simulate action potential propagation in a myelinated axon.
\end{enumerate}

\begin{thebibliography}{7}
    \bibitem{dayan2001} P Dayan, LF Abbott, ``Theoretical neuroscience'',
        MIT Press 2001
\end{thebibliography}

\vfill
\centerline{\CAP Contact}
\CAP

\begin{tabular}{lll}
\underline{Urs Bergmann} & Phone: 2093-8924 & Email:
urs.bergmann(at)biologie.hu-berlin.de \\
Matthias Guggenmos & & Email: matthias.guggenmos(at)bccn-berlin.de \\
Richard Kempter \hfill & Phone: 2093-8925 \hfill & Email:
r.kempter(at)biologie.hu-berlin.de \\
\end{tabular}

\end{document}
                        
